\documentclass[../../../../SoftwareRequirementsSpecification.tex]{subfiles}
\begin{document}
% Richiede le macro definite in src/macros/useCaseMacros.tex
% (incluse nel preambolo di SoftwareRequirementsSpecification.tex)

\ucstart
\begin{xltabular}{\textwidth}{|l|X|}
  \hline
  \ucid{PT-16} \\ \hline
  \ucname{Duplicazione Scheda} \\ \hline
  \ucactor{PT} \\ \hline
  \ucdescription{Il PT crea una copia di una scheda che ha creato.
    Il PT può modificare la copia anche se la scheda originale ha
    assegnazioni.} \\ \hline
  \ucprecond{Il PT deve essere autenticato. Deve esistere almeno una
    scheda creata dal PT.} \\ \hline
  \ucpostcond{Il sistema crea una nuova scheda senza assegnazioni.
    La nuova scheda ha gli stessi blocchi mensili, le stesse sessioni
    di allenamento e gli stessi esercizi della scheda originale. La
    scheda originale non cambia.} \\ \hline
  \ucmainflow{
    \item Il PT apre la pagina delle proprie schede.
    \item Il sistema mostra l'elenco delle schede create dal PT.
    \item Il PT preme il pulsante "Duplica" su una scheda
      dell'elenco.
    \item Il sistema mostra il form con il nome della copia,
      precompilato con il nome della scheda originale seguito da
      "(copia)".
    \item Il PT modifica il nome, se vuole, e preme il pulsante
      "Conferma".
    \item Il sistema verifica che il nome non sia già usato da
      un'altra scheda del PT.
    \item Il sistema crea la copia della scheda.
    \item Il sistema reindirizza il PT alla configurazione delle
      sessioni di allenamento della copia (caso d'uso PT-04).
  } \\ \hline
  \ucaltflow{
    \textbf{6a. Nome già presente} \newline
    - Il sistema mostra un messaggio di errore ("Esiste già una
    scheda con questo nome"). \newline
    - Il flusso riprende dal passo 4.
  } \\ \hline
  \ucnote{
    La copia non ha assegnazioni, quindi il PT può modificarne i
    blocchi mensili, le sessioni di allenamento e gli esercizi con il
    caso d'uso PT-04. Le assegnazioni della scheda originale non
    cambiano: gli atleti che la seguono continuano a vedere la scheda
    originale (caso d'uso A-03).
  } \\ \hline
\end{xltabular}
\end{document}
